The predeclared screen used MORL-Baselines 1.3.0 Envelope and the pinned
upstream Minecart example, with the repository's frozen $\gamma=0.99$ and reward
scales. Each of three seeds ran exactly 100,000 training interactions and
99,901 gradient updates. Gaussian behavior/update preferences were retained;
this is not a fixed-pool or compute-matched allocation comparator. Validation
used all 32 course preferences, three episodes each, at fixed steps 50,000 and
100,000. Training RNGs and environments were isolated from validation. The
original CUDA smoke exposed a replay-index transport bug; a tested host-index
bridge preserved the sampled indices/priorities and allowed the pinned learner
to run. No learning rule or hyperparameter was tuned on these validation results.

\begin{table}[htbp]
\centering\small
\caption{Envelope feasibility: positive ore episodes out of 96 fixed cases at each checkpoint.}
\begin{tabular}{rrrrrr}
\toprule
Seed & Steps & Ore 1 & Ore 2 & HV & Expected utility\\
\midrule
20261121 & 50,000 & 0 & 0 & 0.001477 & -0.5444\\
20261121 & 100,000 & 0 & 0 & 0.001448 & -0.5512\\
20261122 & 50,000 & 0 & 0 & 0.001421 & -0.5868\\
20261122 & 100,000 & 2 & 2 & 0.001359 & -0.7790\\
20261123 & 50,000 & 2 & 2 & 0.001283 & -0.6058\\
20261123 & 100,000 & 0 & 0 & 0.001420 & -0.5967\\
\bottomrule\end{tabular}\end{table}

The completed cohort used 300,000 training interactions, 299,703 updates and 530,664 separate validation environment steps across 576 episodes. Final-test steps were zero. At least one final checkpoint expressed positive ore in some validation cases. The per-seed counts above delimit that observation; they do not establish robust task coverage or an adaptive-allocation benefit. All model files and validation tensors passed hash, shape, finite-value, budget and metric checks. This numerical audit is not an independent semantic review. Neither outcome unlocks the failed matrix-v1/F1 or final-test workflow.
